\subsection{17. Updated Virophage Taxonomy and Distinction from Polinton-like Viruses}
\textbf{OSTI:} \texttt{2475938}\quad\textbf{Authors:} Roux et al. (2023)\quad\textbf{Domain:} Virology / bioinformatics\\
\scorebadge{Coverage}{8}\quad\scorebadge{Agreement}{10}\quad\textbf{Total:} 18/20\par\smallskip
\noindent\textbf{Coverage rationale.} \textbf{Wave~2 scale-up:} 279-genome NCBI bulk-fetch ($21\times$ the 13-genome baseline); 4,676 predicted proteins; 70-genome partitioned 4-marker concatenated ML tree (1,403~aa, 4 LG$+$G partitions, IQ-TREE 2 with 1000 UFBoot). Classification: 75 virophage, 113 PLV/Polinton-like, 5 partial, 86 unclassified. Mavirus/Sputnik/Aquatic-Lavidaviridae clades $+$ Maveriviricetes/Polintoviricetes boundary recovered matching paper's Fig.~3. Residual gap: IMG/VR uncultivated tail (JGI Globus auth required), AAI/Mash/MCL clustering, CheckV filtering.\par
\noindent\textbf{Agreement rationale.} Family-level backbone congruent with paper: Mavirus clade (NC\_015230, marine assemblies), Sputnik clade (UFBoot $\geq 94$), Aquatic-virophage assemblage, Ace-Lake deep branch. HMM partition perfect on curated subset. Scale from 13 to 279 genomes did not disrupt topology.\par\smallskip
\noindent\textbf{What reproduced:}
\begin{itemize}[leftmargin=1.4em,itemsep=0.2em,topsep=0.2em]
\item Prodigal ORF prediction
\item HMMER3 scan against paper's 19 virophage + 1 PLV HMMs
\item Clean virophage vs PLV/adintovirus separation
\item Four-marker MAFFT + trimAl alignments
\item IQ-TREE ML phylogenies with 1000 UFBoot
\item Backbone topology: Mavirus basal, Sputnik early, SW01+YSLV4 sister
\item HMM heatmap and 4-marker congruence figures
\end{itemize}\noindent\textbf{What missing:}
\begin{itemize}[leftmargin=1.4em,itemsep=0.2em,topsep=0.2em]
\item Full paper-scale genome collection (hundreds of genomes)
\item ICTV-scale family/genus assignments
\item GC-content and gene-content discriminant analysis
\item Deep-clade (Preplasmiviricetes) re-classification
\item Synteny / gene-order analysis across virophages
\item Metagenomic-contig based new virophage discovery
\end{itemize}\noindent\textbf{Follow-on research questions:}
\begin{enumerate}[leftmargin=1.6em,itemsep=0.2em,topsep=0.2em,label=Q\arabic*.]
\item Does adding the 32 newly published metagenome-assembled virophage genomes (2023-2024 IMG/VR release) shift the Mavirus-basal topology, or are the paper's conclusions robust
\item Can a protein language model (ESM-2 or ProtTrans) embedding of MCP and ATPase sequences distinguish virophages from PLVs/adintoviruses without HMM profiles, and at what sequence-identity floor does HMMER fail while LM embeddings still succeed
\item What is the minimum marker-gene subset (MCP alone vs 2/3/4 markers) needed to recover the paper's family-level topology at \textgreater{}= 95 UFBoot support
\item Under Bayesian concordance factor analysis, do the four marker trees support a single evolutionary history or reveal reticulation (horizontal transfer) at the Sputnik/Mavirus split
\item When reference-free clustering (vConTACT3 or mash-based) is applied to the paper's genome set, does it produce the same family boundaries, or does it over-/under-split relative to the HMM + phylogeny approach
\end{enumerate}

